\section{Prediction from short histories}\label{sec:prediction}

The incidence reduction supplies a probability law on point--center pairs whose
point endpoints in a random wedge have mutual information \(D=o(\log t)\).
We now turn this information bound into predictions from short histories of
point tests.  The prediction uses only the history of the point; its accuracy
is measured against a test that may depend on the center.  We first prove this
statement for an arbitrary finite probability law and then give the two forms
needed for coordinate quantization and coordinate equality.

The proof sums the information contributed by successive tests. Since
that sum is at most \(D\), a sufficiently long history leaves little
new dependence on average. Related short conditioning arguments occur in
Ahlswede's
wringing method \cite[Section~4, Lemma~3]{Ahlswede1982}.
Here the proof must additionally compare each endpoint with its own
posterior law and leave the next test independent of the selected history.

\subsection{The probability laws and the entropy comparison}

Let \(P,Q\) be finite nonempty sets and let \(\omega(p,q)\) be a probability
mass function on \(P\times Q\), with marginals \(\mu\) and \(\kappa\).
Delete atoms of zero marginal mass.  Write
\[
 \omega(p\mid q)=\frac{\omega(p,q)}{\kappa(q)},
 \qquad
 \omega(q\mid p)=\frac{\omega(p,q)}{\mu(p)}.
\]
The associated wedge law and its point-pair marginal are
\begin{equation}\label{pred:wedge-law}
 \Omega(p,q,r)=\kappa(q)\omega(p\mid q)\omega(r\mid q),
 \qquad
 \pi(p,r)=\sum_{q\in Q}\Omega(p,q,r).
\end{equation}
Both marginals of \(\pi\) are \(\mu\). Relative entropy is defined in
\eqref{inc:relative-entropy}. For finite probability laws \(a,b\), set
\[
 \TV(a,b)=\frac12\sum_x\abs{a(x)-b(x)}.
\]
The information available to the argument is
\begin{equation}\label{pred:information}
 D=\KL(\pi\Vert\mu\otimes\mu).
\end{equation}
It is finite, since \(\pi(p,r)>0\) implies \(\mu(p)\mu(r)>0\).

We use the standard chain rule for relative entropy
\cite[Theorem~2.5.3]{CoverThomas2006}:
\begin{equation}\label{pred:chain-rule}
 \KL(a_{XY}\Vert b_{XY})
 =\KL(a_X\Vert b_X)
   +\sum_x a_X(x)\KL(a_{Y\mid x}\Vert b_{Y\mid x}).
\end{equation}
Nonnegativity and the chain rule imply that applying a deterministic map
cannot increase relative entropy: adjoin the map to the outcome and then
discard the nonnegative conditional term. We also use Pinsker's inequality
\cite[Lemma~11.6.1]{CoverThomas2006}, which in our natural-logarithm
convention is
\begin{equation}\label{pred:pinsker}
 2\TV(a,b)^2\le\KL(a\Vert b).
\end{equation}
Finally, total variation controls bounded expectations by
\begin{equation}\label{pred:tv-expectation}
 \abs{\E_a F-\E_b F}\le2\norm{F}_{\infty}\TV(a,b).
\end{equation}

A test \(T\) is a random element of a probability space, independent of
\((p,q,r)\), which specifies a code map \(c_T:P\to[M]\), where
\([M]=\{1,\ldots,M\}\).  The test space may be infinite; all maps and
integrands below are assumed measurable.  Only finite probability laws of
codes are used.  All vector norms and inner products in this section are
Euclidean.  For a fixed history map \(H:P\to\mathcal H\), put
\begin{equation}\label{pred:posterior-point}
 \mu_h(x)=\frac{\mu(x)\mathbf1_{\{H(x)=h\}}}{\mu(H=h)}
 \quad\text{when }\mu(H=h)>0.
\end{equation}
Its values on histories of zero mass are immaterial.

\begin{lemma}[Prediction device]\label{pred:device}
Let \(\omega,\pi,D\) be as in \eqref{pred:wedge-law}--\eqref{pred:information},
let \(K,M\ge1\) be integers, and let \(T_1,\ldots,T_K\) be independent tests
with the same law as \(T\).  Each test specifies, in addition to \(c_T\),
a feature map \(f_T:[M]\to\R^d\) and a coefficient map
\(g_T:Q\to\R^d\), for an integer \(d\ge1\) and constants
\(B,B'\ge0\), satisfying
\[
 \norm{f_T(c)}\le B,\qquad \norm{g_T(q)}\le B'
\]
for all codes, centers and tests.  Choose \(J\) uniformly from \([K]\),
independently of all tests, and retain only the strict prefix
\[
 H(p)=(c_{T_1}(p),\ldots,c_{T_{J-1}}(p)).
\]
Conditioned on \(J,T_1,\ldots,T_{J-1}\), let \(T\) be a fresh test and define
\[
 m_{T,h}=\sum_{x\in P}\mu_h(x)f_T(c_T(x)).
\]
Then
\begin{equation}\label{pred:device-bound}
 \E_{J,T_1,\ldots,T_{J-1},T}
 \left|\E_{(p,q)\sim\omega}
 \ip{f_T(c_T(p))-m_{T,H(p)}}{g_T(q)}\right|
 \le4BB'\left(\frac DK\right)^{1/4}.
\end{equation}
Every realized history has at most \(M^{K-1}\) values.  In particular, a
history can be fixed before drawing the fresh test so that its average
absolute discrepancy in \eqref{pred:device-bound} is at most the right side.
\end{lemma}

\begin{proof}
First freeze all \(K\) tests.  Apply their complete transcript map to both
coordinates of the pair \((p,r)\).  The actual transcript law is the image of
\(\pi\); the reference law is the image of \(\mu\otimes\mu\).
Data processing bounds their relative entropy by \(D\).
For \(j\in[K]\), write \(H_j\) for the transcript strictly before step \(j\).
Conditional on prefixes \((h,h')\), the reference next-code law is exactly
\[
 b_{j,h,h'}=
 (c_{T_j})_*\mu_h\ \otimes\ (c_{T_j})_*\mu_{h'}.
\]
Here \(\mu_h\) is formed with \(H_j\).  Let \(a_{j,h,h'}\) denote the actual
conditional law of the next two codes.  All prefix pairs of positive actual
probability have positive reference probability.  Iterating
\eqref{pred:chain-rule} therefore gives
\begin{equation}\label{pred:conditional-kl}
 \sum_{j=1}^K
 \E_{(p,r)\sim\pi}
 \KL\bigl(a_{j,H_j(p),H_j(r)}\Vert
          b_{j,H_j(p),H_j(r)}\bigr)\le D.
\end{equation}
Consequently, by \eqref{pred:pinsker} and Cauchy--Schwarz,
\begin{equation}\label{pred:conditional-tv}
 \frac1K\sum_{j=1}^K
 \E_{(p,r)\sim\pi}
 \TV\bigl(a_{j,H_j(p),H_j(r)},b_{j,H_j(p),H_j(r)}\bigr)
 \le\sqrt{\frac{D}{2K}}.
\end{equation}

For this fixed test sequence and step \(j\), subtract the prediction belonging
to each point's own prefix:
\[
 U_j(p)=f_{T_j}(c_{T_j}(p))-
            \sum_x\mu_{H_j(p)}(x)f_{T_j}(c_{T_j}(x)).
\]
The norm is at most \(2B\).  Under each reference conditional product, both
residuals are centered and independent, so the conditional expectation of
\(\ip{U_j(p)}{U_j(r)}\) is zero.  This inner product has absolute value at
most \(4B^2\).  Thus \eqref{pred:tv-expectation} and
\eqref{pred:conditional-tv} imply
\begin{equation}\label{pred:residual-variance}
 \frac1K\sum_{j=1}^K
 \E_{(p,r)\sim\pi}\ip{U_j(p)}{U_j(r)}
 \le8B^2\sqrt{\frac{D}{2K}}.
\end{equation}
Each summand is nonnegative, for conditional independence in the wedge law
gives the exact identity
\begin{equation}\label{pred:wedge-square}
 \E_{(p,r)\sim\pi}\ip{U_j(p)}{U_j(r)}
 =\E_{q\sim\kappa}
       \left\|\E_{p\sim\omega(\cdot\mid q)}U_j(p)\right\|^2.
\end{equation}
It follows, by Cauchy--Schwarz first over \(q\) and then over the tests and
step, that the average absolute discrepancy is at most
\[
 B'\left(8B^2\sqrt{\frac{D}{2K}}\right)^{1/2}
 \le4BB'(D/K)^{1/4}.
\]
Everything proved so far holds after averaging the initially frozen tests.
The residual at step \(j\) uses only \(T_1,\ldots,T_j\).  Since the tests
are independent and identically distributed, its joint law with the strict
prefix is unchanged when \(T_j\) is replaced by a fresh test after that
prefix.  This proves \eqref{pred:device-bound}.  Finally average over the
strict prefixes and choose one attaining at most the average.  There are at
most \(M^{j-1}\le M^{K-1}\) possible histories for a prefix of length \(j-1\).
\end{proof}

The important comparison in this proof is conditional on two possibly
different histories.  The reference law uses the first endpoint's prediction
at the first coordinate and the second endpoint's prediction at the second.
No independence of the actual endpoints after their histories have been
specified is asserted.  Their independence is used only conditional on the
center, in \eqref{pred:wedge-square}.

\subsection{Acceptance tests and the prediction coupling}

\begin{corollary}[Finite-code acceptance]\label{pred:acceptance}
In the setting of \Cref{pred:device}, suppose each test specifies a set
\(\mathcal A_T(q)\subseteq[M]\) of accepted codes for every \(q\in Q\).
Put
\[
 s_0=\E_T\Prob_{(p,q)\sim\omega}
                \{c_T(p)\notin\mathcal A_T(q)\}.
\]
There is a fixed strict-prefix history \(H\), with at most \(M^{K-1}\)
values, such that
\begin{equation}\label{pred:acceptance-bound}
 \E_T\sum_{p,q}\omega(p,q)
       \sum_x\mu_{H(p)}(x)
             \mathbf1_{\{c_T(x)\notin\mathcal A_T(q)\}}
 \le s_0+4\sqrt M(D/K)^{1/4}.
\end{equation}
The fresh test \(T\) is independent of the fixed history construction.
\end{corollary}

\begin{proof}
Use \(f_T(c)=e_c\in\R^M\), of norm one, and use the acceptance indicator
vector as \(g_T(q)\), of norm at most \(\sqrt M\).  Actual and predicted
failure probabilities differ by the negative of their acceptance
probabilities.  Averaging \eqref{pred:device-bound} and then fixing a strict
prefix proves the assertion.
\end{proof}

We record the sampling laws used with this result.  For any fixed history
\(H\) and integer \(r\ge1\), draw one point, one prediction, and \(r\)
neighbors according to
\begin{equation}\label{pred:shared-prediction-law}
 \Prob\{p,\widehat p,q_1,\ldots,q_r\}
 =\mu(p)\mu_{H(p)}(\widehat p)
             \prod_{j=1}^r\omega(q_j\mid p),
\end{equation}
and draw the fresh test independently of this entire tuple.  The two point
marginals are both \(\mu\), each \((p,q_j)\) has law \(\omega\), and each
\(q_j\) has marginal \(\kappa\).  In particular each pair
\((p,\widehat p)\) is independent of the fresh test, although its two points
need not be independent.  These assertions follow by summing
\eqref{pred:shared-prediction-law}, using
\(\sum_{p:H(p)=h}\mu(p)\mu_h(x)=\mu(x)\mathbf1_{\{H(x)=h\}}\).
For every fixed history,
\begin{equation}\label{pred:point-collision}
 \Prob\{p=\widehat p\}
 =\sum_h\frac{\sum_{p:H(p)=h}\mu(p)^2}{\mu(H=h)}
 \le\abs{H(P)}\max_p\mu(p).
\end{equation}
Given \(H(p)\), the prediction is independent of all neighbors in
\eqref{pred:shared-prediction-law}.  Applying a union bound to
\eqref{pred:acceptance-bound} therefore bounds the chance that at least one
of the \(r\) acceptance comparisons rejects this \emph{same} predicted point
at the same fresh test by
\(r(s_0+4\sqrt M(D/K)^{1/4})\).  The union bound needs no independence of
the rejection events.

There is also a useful equivalent way to draw the point pair: draw
\(h\) with probability \(\mu(H=h)\), then draw two points independently
from \(\mu_h\).  This again has both marginals \(\mu\), and its collision
probability satisfies \eqref{pred:point-collision}.  Once \(H\) is fixed,
this pair law can be used unchanged for every fresh test. Sampling a second
object from the first object's fiber also appears in the entropy argument
of Katz and Tardos \cite[Section~2, proof of Lemma~1]{KatzTardos}; the
identities above specify the marginals and independent test needed here.

For the coordinate-quantization application, a test will be an independently
chosen field embedding, a code will specify bounded coordinate bins together
with an overflow code, and \(\mathcal A_T(q)\) will consist of codes achieved
by actual neighbors of \(q\).  The exact hypotheses needed from that
construction are a fixed finite alphabet of size \(M\) and acceptance of
every actual edge, so \(s_0=0\).  If \(K=o(\log t)\),
\(D/K=o(1)\), and \(\max_p\mu(p)\le t^{-1+o(1)}\), then the prediction
failure and the collision probability both tend to zero.  Formula
\eqref{pred:shared-prediction-law} with \(r=3\) is the simultaneous
three-neighbor experiment.  Geometric estimates valid for all couplings with
these marginals remain valid for each fixed history.  More generally,
additional nonnegative error quantities can be averaged together with the
prediction error before choosing a prefix; this always leaves the test fresh.

\subsection{Equality of one coordinate}

The next application allows an arbitrarily large set of state labels without
paying for its size in the feature norm.  A state is a pair in a product
\(\mathcal A\times\mathcal B\).  For a state \(y=(a,b)\), define its
\emph{strict cross}
\begin{equation}\label{pred:strict-cross}
 R_y=\{(a',b'):
           (a'=a\text{ and }b'\ne b)
           \text{ or }(a'\ne a\text{ and }b'=b)\}.
\end{equation}
Thus exactly one coordinate agrees.  A fresh test specifies state maps
\(y_T(x)=(A_T(x),B_{1,T}(x))\) for all \(x\in P\cup Q\).
Only the finitely many labels attained on these vertex sets need be retained.
The labels and their equality relations are assumed measurable in \(T\).

\begin{corollary}[Prediction of strict crosses]\label{pred:cross}
Let \(\omega,\pi,D\) be as above, let \(K\ge1,L\ge2\) be integers, and
let \(T\) be an independent random test giving the state maps just described.
Set
\[
 s_0=\E_T\Prob_{(p,q)\sim\omega}
                       \{y_T(p)\notin R_{y_T(q)}\}.
\]
There exists a fixed history \(H:P\to\mathcal H\), with
\(\abs{H(P)}\le L^{2K}\), for which the laws of the \emph{true}, unhashed
states
\[
 \nu_{h,T}=(y_T)_*\mu_h
\]
satisfy
\begin{equation}\label{pred:cross-bound}
 \E_T\E_{(p,q)\sim\omega}
       [1-\nu_{H(p),T}(R_{y_T(q)})]
 \le s_0+\frac4L+18(D/K)^{1/4}.
\end{equation}
The history is formed from earlier tests and does not use \(T\).
\end{corollary}

\begin{proof}
Augment each test by two fresh independent random maps from its attained
coordinate labels into \([L]\), assigning independent uniform values to
distinct labels.  Numbering attained labels by their first occurrence in a
fixed vertex order makes this construction measurable.  Use the same map
on both vertex sets for each coordinate,
and use independent maps for the two coordinates and for different tests.
The point code is its two hashed coordinates, so the alphabet has size
\(M=L^2\).  For a code \((a,b)\in[L]^2\), use the feature
\[
 f(a,b)=(e_a,e_b,e_{(a,b)})\in\R^L\oplus\R^L\oplus\R^{L^2},
 \qquad \norm{f(a,b)}=\sqrt3.
\]
For the hashed center state \((a_q,b_q)\), use
\[
 g(q)=(e_{a_q},e_{b_q},-2e_{(a_q,b_q)}),
 \qquad \norm{g(q)}=\sqrt6.
\]
Their inner product is
\(\mathbf1_{\{a=a_q\}}+\mathbf1_{\{b=b_q\}}
-2\mathbf1_{\{(a,b)=(a_q,b_q)\}}\), exactly the indicator of the hashed
strict-cross relation.  Therefore \Cref{pred:device} compares its actual and
predicted probabilities with error at most
\(4\sqrt{18}(D/K)^{1/4}\), averaged over strict prefixes and the fresh
augmented test.

For any two fixed true states, hashing preserves an equality, and creates
an equality between two unequal coordinate labels with probability \(1/L\).
Thus a true strict-cross indicator differs from its hashed indicator with
probability at most \(2/L\).  This applies both to the actual state pair and
to the predicted pair: conditional on the earlier prefix, the fresh true
test and the chosen vertices, the current hashes are still independent
uniform maps.  In particular the predicted point is sampled by
\eqref{pred:posterior-point} using no current hash.  Comparing actual true,
actual hashed, predicted hashed, and predicted true failures gives, in that
order, total extra error at most
\(4/L+4\sqrt{18}(D/K)^{1/4}\).  Average the current test and current hashes
before fixing a prefix, and use \(4\sqrt{18}<18\).  The resulting true-state
expression no longer involves the current hashes.  Its history has at most
\((L^2)^{K-1}\le L^{2K}\) values.
\end{proof}

One form of \Cref{pred:cross} that will be convenient is an integrated
version.  Let \(\mathcal T\) be a measurable family of state tests with finite
measure \(W\), and let \(m(\tau)\) bound actual strict-cross failure at
\(\tau\).  For \(W>0\), sample \(T\) from the normalized measure
\(\dd\tau/W\).  The resulting fixed history satisfies
\begin{equation}\label{pred:integrated-cross}
 \int_{\mathcal T}\E_{(p,q)\sim\omega}
 [1-\nu_{H(p),\tau}(R_{y_\tau(q)})]\dd\tau
 \le\int_{\mathcal T}m(\tau)\dd\tau
       +W\bigl(4/L+18(D/K)^{1/4}\bigr).
\end{equation}
For \(W=0\), the assertion holds with the empty history.  In the later level
application, \(\mathcal T\) is the common-level set with its weighted level
measure, and \(m\) is the incidence inconsistency probability.  Consistency
implies the strict-cross relation, which is the only geometric hypothesis
needed for \eqref{pred:integrated-cross}.  We suppress the test subscript and
write \(\nu_h\) when the level being tested is fixed.

\subsection{Choosing the history length}

There is no conflict between small prediction error and a subpolynomial
number of histories.  Write \(T_0=\log t\), and suppose \(D=o(T_0)\).
For example, take
\[
 K=\left\lceil\sqrt{(D+1)T_0}\right\rceil.
\]
Then \(K\to\infty\), \(K=o(T_0)\), and \(D/K\to0\).
For a fixed alphabet this gives \(M^K=t^{o(1)}\).
For hashing one may additionally take
\[
 L=\left\lfloor\exp\sqrt{T_0/K}\right\rfloor.
\]
Eventually \(L\ge2\), \(L\to\infty\), and
\(K\log L=o(T_0)\), so \(L^{2K}=t^{o(1)}\).
These statements require no quantitative rate for \(D/T_0\).
If a later argument requires an integer \(k\to\infty\) satisfying
\(kK\log L=o(T_0)\) and \(k^{20}\eta_t\to0\) for an additional positive
sequence \(\eta_t\to0\), it suffices to choose
\[
 k=\left\lfloor\min\left\{
       \left(\frac{T_0}{K\log L}\right)^{1/2},
       \eta_t^{-1/40}\right\}\right\rfloor.
\]
Further finitely many slow-growth requirements can be imposed by taking a
smaller such diverging sequence.  Thus the prediction lemma supplies both
vanishing error and the small history counts required in the subsequent
geometric arguments.
